\newpage

\section{Пункт А: Дифракция Френеля на щели}  
\subsection{Описание экспериментальной установки}  
Экспериментальная установка для наблюдения дифракции Френеля включала ртутную лампу, излучение которой проходило через оранжевый светофильтр (\(\lambda = 578 \, \text{нм}\)), формируя монохроматический пучок. Коллиматор, состоящий из щели \(S_1\) и собирающей линзы \(O_1\), преобразовывал свет в параллельный пучок, который освещал исследуемую щель \(S_2\). Ширина щели \(S_2\) (\(b\)) регулировалась с высокой точностью, а дифракционная картина наблюдалась в плоскости \(П\) с помощью микроскопа \(M\), сфокусированного на эту плоскость.  

Ключевым параметром эксперимента являлось расстояние \(z\) от щели \(S_2\) до плоскости наблюдения. При изменении \(z\) варьировалось число открытых зон Френеля (\(m_{\text{max}} = \frac{b^2}{4\lambda z}\)), что влияло на структуру дифракционной картины. Например, при \(z \to 0\) (\(m \gg 1\)) наблюдался переход к геометрической оптике, тогда как при \(z \to \infty\) (\(m \ll 1\)) картина переходила в режим Фраунгофера.  

Микроскоп позволял фиксировать положение минимумов и максимумов, соответствующих чётному и нечётному числу открытых зон. Для анализа использовались границы зон Шустера (\(\xi_m = \pm\sqrt{m z \lambda}\)), определявшие распределение интенсивности в плоскости \(П\).  

\begin{figure}[H]  
\centering  
\includegraphics[width=1\textwidth]{fresnel_setup}  
\caption{Схема установки для наблюдения дифракции Френеля}  
\label{fig:fresnel_setup}  
\end{figure}  

При фиксированной ширине щели \(b\) и изменении расстояния \(z\) можно наблюдать, как число открытых зон Френеля \(m_{\text{max}} = \frac{b^2}{4\lambda z}\) влияет на структуру дифракционной картины. При открытии \(m\) зон Шустера на экране возникает \(m-1\) тёмных полос. 

\subsection{Обработка результатов измерений}  
Экспериментальные данные получены при ширине щели \( b = 0.30 \pm 0.05 \, \text{мм} \). Координаты микроскопа \( z \) при наблюдении \( n \) тёмных полос приведены в таблице \ref{tab:fresnel_data} (погрешность измерения по линейке 5 \%). Число открытых зон Френеля рассчитывалось как \( m = n + 1 \). Щель была чётко видна на 57.5 см, будем считать это расстояние за точку отсчёта $z = 0$.

% Please add the following required packages to your document preamble:
% \usepackage{longtable}
% Note: It may be necessary to compile the document several times to get a multi-page table to line up properly
\begin{longtable}[c]{|c|c|c|c|}
\caption{Зависимость расстояния \( z \) от числа открытых зон \( m \)}
\label{tab:fresnel_data}\\
\hline
Полосы, n & Зоны, m & Координата по линейке, см & $z$, см \\ \hline
\endfirsthead
%
\endhead
%
1         & 2       & 55.35                     & 2.15    \\ \hline
2         & 3       & 56.1                      & 1.4     \\ \hline
3         & 4       & 56.45                     & 1.05    \\ \hline
4         & 5       & 56.7                      & 0.8     \\ \hline
5         & 6       & 56.85                     & 0.65    \\ \hline
\end{longtable}

\subsubsection*{Построение графика и определение ширины щели}  
Для проверки теоретического соотношения \(4\lambda m_{\text{max}} = \frac{b^2}{z}\) построен график зависимости \(4\lambda m\) от \(1/z\) (рис. \ref{fig:4lambda_vs_invz}). Экспериментальные данные аппроксимированы прямой линией методом наименьших квадратов:  
\[
y = (8.78 \pm 0.14) \times 10^{-10} \ \cdot x.
\]  
Наклон прямой \(k = (8.78 \pm 0.14) \times 10^{-10} \, \text{м}^2\) соответствует квадрату ширины щели:  
\[
b = \sqrt{k} = \sqrt{8.78 \times 10^{-10}} = (0.296 \pm 0.024) \, \text{мм}.
\]  
Это значение согласуется с измеренным \(b_{\text{эксп}} = 0.3 \, \text{мм}\) в пределах погрешности (\(\varepsilon = 1\%\)).

\begin{figure}[h]
\centering
\includegraphics[width=0.8\textwidth]{4lambda_vs_invz}
\caption{Зависимость \(4\lambda m\) от \(1/z\). Наклон прямой равен \(b^2\)}
\label{fig:4lambda_vs_invz}
\end{figure}

\subsubsection*{Проверка критерием \(\chi^2\)}  
Для оценки согласия эксперимента с теорией вычислен критерий \(\chi^2\):  
\[
\chi^2 = \sum \frac{(y_i - k \cdot x_i)^2}{\sigma_{y_i}^2} = 2.1 \, \, (\nu = 3).
\]  
Критическое значение \(\chi^2_{\text{крит}} = 7.815\) (\(\nu = 3, \alpha = 0.05\)). Так как \(\chi^2 < \chi^2_{\text{крит}}\), гипотеза о линейности \textbf{принимается}.

\subsection{Качественный анализ интенсивности}  
\begin{figure}[H]
\centering
\includegraphics[width=0.6\textwidth]{intensity_distribution}
\caption{Распределение интенсивности при дифракции на краю экрана}
\label{fig:intensity}
\end{figure}

На границе света и тени (рис. \ref{fig:intensity}) интенсивность плавно спадает от максимума к нулю. Спираль Корню (рис. \ref{fig:cornu}) иллюстрирует, как амплитуда результирующего колебания зависит от числа зон Френеля:  
- Вблизи края щели (\(m \sim 1\)) наблюдается чередование ярких и тёмных полос.  
- Вдали от края (\(m \gg 1\)) интерференционные эффекты затухают.

\begin{figure}[H]
\centering
\includegraphics[width=0.5\textwidth]{cornu_spiral}
\caption{Спираль Корню для расчёта амплитуды при дифракции}
\label{fig:cornu}
\end{figure}

\subsection{Связь числа полос с шириной щели}  
Число наблюдаемых тёмных полос \(n\) определяется числом зон Френеля \(m = n + 1\), укладывающихся в ширину щели:  
\[
m = \frac{b^2}{4\lambda z}.
\]  
При фиксированной \(b\) увеличение \(z\) уменьшает \(m\), что приводит к уменьшению числа полос. Например, для \(b = 0.3 \, \text{мм}\) и \(z = 10 \, \text{см}\):  
\[
m = \frac{(0.3 \times 10^{-3})^2}{4 \cdot 546.1 \times 10^{-9} \cdot 0.01} \approx 4.1 \, \, \Rightarrow \, n = 3 \, \text{полосы}.
\]  
Экспериментально наблюдалось 3 полосы.

Ясно, что при уменьшении щели будет уменьшаться число видимых полос, согласно формуле выше.

\section{Пункт Б: Дифракция Фраунгофера на щели}  
\subsection{Теоретическое описание явления}  
Дифракция Фраунгофера возникает при выполнении условия \( b \ll \sqrt{\lambda z} \), когда ширина щели \( b \) значительно меньше размера первой зоны Френеля. В этом случае дифракционная картина формируется в дальней зоне (условно — на бесконечности) и характеризуется угловым распределением интенсивности, не зависящим от расстояния \( z \). Для наблюдения явления в лабораторных условиях используется дополнительная линза \( O_2 \), которая фокусирует параллельные лучи в своей фокальной плоскости, имитируя бесконечно удалённый экран (рис. \ref{fig:fraunhofer_scheme}).

\begin{figure}[h]  
\centering  
\includegraphics[width=1\textwidth]{fraunhofer_scheme}  
\caption{Схема установки для наблюдения дифракции Фраунгофера}  
\label{fig:fraunhofer_scheme}  
\end{figure}  

Установка включает: 

- Ртутную лампу с зелёным светофильтром (\( \lambda = 546.1 \, \text{нм} \)).  

- Коллиматор (\( O_1 \)), формирующий параллельный пучок света.  

- Щель \( S_2 \) с регулируемой шириной \( b \).  

- Объектив \( O_2 \) (\( f_2 = 15 \, \text{см} \)), создающий изображение в фокальной плоскости.  

- Микроскоп для наблюдения дифракционной картины.  

\subsubsection*{Ключевые соотношения}  
1. \textbf{Условие минимумов интенсивности}:  
   Угловое положение \( n \)-го минимума определяется формулой:  
   \[
   \sin\Theta_n = n \frac{\lambda}{b}, \quad n = \pm 1, \pm 2, \ldots  
   \]  
   Для малых углов (\( \Theta_n \ll 1 \)) это упрощается до:  
   \[
   \Theta_n \approx n \frac{\lambda}{b}.  
   \]  

2. \textbf{Линейное положение минимумов}:  
   В фокальной плоскости линзы \( O_2 \) координата \( X_n \) минимума связана с углом \( \Theta_n \):  
   \[
   X_n = f_2 \cdot \Theta_n = n \frac{\lambda f_2}{b},  
   \]  
   где \( f_2 \) — фокусное расстояние линзы.  

3. \textbf{Ширина центрального максимума}:  
   Расстояние между первыми минимумами (\( n = \pm 1 \)):  
   \[
   \Delta X = 2 \frac{\lambda f_2}{b}.  
   \]  

\subsection{Обработка результатов измерений}

\begin{figure}[h]  
\centering  
\includegraphics[width=1\textwidth]{diffraction_pattern}  
\caption{Дифракция Фраунгофера на одной щели}  
\label{fig:diffraction_pattern}  
\end{figure}  

По данным окулярной шкалы микроскопа определены координаты минимумов \( x_m \). Цена деления: \( 0.02 \, \text{мм} \). Погрешности: \( \pm 1 \, \text{дел} \) для \( m = 1–3 \), \( \pm 2 \, \text{дел} \) для \( m = 4–6 \).

% Добавьте в преамбулу:
% \usepackage{longtable}
\begin{longtable}[c]{|c|c|c|c|}
\caption{Координаты минимумов дифракции Фраунгофера}
\label{tab:fraunhofer_data}\\
\hline
Порядок, $m$ & Деления & $x_m$, мм & $\sigma_{x_m}$, мм \\ \hline 
\endfirsthead

\multicolumn{4}{c}{{\tablename} \thetable\ - продолжение} \\
\hline
Порядок, $m$ & Деления & $x_m$, мм & $\sigma_{x_m}$, мм \\ \hline 
\endhead

1 & 9 & 0.18 & 0.02 \\ \hline
2 & 18 & 0.36 & 0.02 \\ \hline
3 & 27 & 0.54 & 0.02 \\ \hline
4 & 35 & 0.70 & 0.04 \\ \hline
5 & 45 & 0.90 & 0.04 \\ \hline
6 & 53 & 1.06 & 0.04 \\ \hline
\end{longtable}
 
Для определения ширины щели $b$ использовалась зависимость $n \lambda f_2$ от $x_n$, где:  
- $n$ — порядок минимума,  
- $\lambda = 546.1$ нм — длина волны,  
- $f_2 = 11.5$ см = 0.115 м — фокусное расстояние линзы,  
- $x_n$ — координата минимума в мм.  

\begin{longtable}[c]{|c|c|c|c|}
\caption{Данные для построения графика $n\lambda f_2$ от $x_n$}
\label{tab:fraunhofer_calculation} \\
\hline
$n$ & $x_n$, мм & $n\lambda f_2$, м$^2$ & $\sigma_{n\lambda f_2}$, м$^2$ \\ \hline 
\endfirsthead

\multicolumn{4}{c}{{\tablename} \thetable\ — продолжение} \\
\hline
$n$ & $x_n$, мм & $n\lambda f_2$, м$^2$ & $\sigma_{n\lambda f_2}$, м$^2$ \\ \hline 
\endhead

1 & 0.18 & $6.28 \times 10^{-8}$ & $3.14 \times 10^{-9}$ \\ \hline
2 & 0.36 & $1.26 \times 10^{-7}$ & $6.28 \times 10^{-9}$ \\ \hline
3 & 0.54 & $1.88 \times 10^{-7}$ & $9.42 \times 10^{-9}$ \\ \hline
4 & 0.70 & $2.51 \times 10^{-7}$ & $1.26 \times 10^{-8}$ \\ \hline
5 & 0.90 & $3.14 \times 10^{-7}$ & $1.57 \times 10^{-8}$ \\ \hline
6 & 1.06 & $3.77 \times 10^{-7}$ & $1.88 \times 10^{-8}$ \\ \hline
\end{longtable}

\begin{figure}[h]
\centering
\includegraphics[width=1\textwidth]{fraunhofer_plot}
\caption{Зависимость $n\lambda f_2$ от $x_n$. Наклон прямой равен $1/b$}
\label{fig:fraunhofer_plot}
\end{figure}


Метод наименьших квадратов (МНК) дал уравнение прямой:  
\[
y = (3.53 \pm 0.02) \times 10^{-7} \, \text{м}^2/\text{мм} \cdot x,
\]  
где коэффициент наклона $k = 3.53 \times 10^{-7} \, \text{м}^2/\text{мм}$ соответствует ширине щели $b$. 

Итоговый результат:  
\[
b = \left( 0.353 \pm 0.002 \right), \text{мм} \ (\varepsilon = 0.57 \%) .
\]  



\subsubsection*{Сравнение с прямым измерением}  
Ширина щели, измеренная микрометрическим винтом:  
\[
b_{\text{эксп}} = 0.35 \, \text{мм}.
\]  

Так, результат совпадает с теоретическим в пределах погрешности.

\section{Пункт В: Дифракция Фраунгофера на двух щелях}  

\subsection{Теоретическое описание явления}  
Дифракция Фраунгофера на двух щелях представляет собой интерференцию волн, прошедших через две параллельные узкие щели. Схема установки для наблюдения этого явления приведена на рис. \ref{fig:double_slit_setup}. В отличие от предыдущих экспериментов, щель $S_2$ заменена на экран $\Theta$ с двумя щелями. Ширина каждой щели $D$ значительно меньше расстояния между ними ($D \ll d$). Линза $O_2$ фокусирует дифракционную картину в своей фокальной плоскости, где наблюдаются интерференционные полосы.  

\begin{figure}[h]  
\centering  
\includegraphics[width=1\textwidth]{double_slit_setup}  
\caption{Схема установки для наблюдения дифракции Фраунгофера на двух щелях}
\label{fig:double_slit_setup}  
\end{figure}  

Интенсивность света в плоскости наблюдения определяется суперпозицией волн от двух щелей. Каждая щель создаёт дифракционную картину, характерную для одиночной щели (штриховая линия на рис. \ref{fig:double_slit_setup}), а интерференция между волнами приводит к появлению узких дополнительных полос. Угловое положение $n$-го интерференционного максимума задаётся условием:  
\[
d \cdot \sin \theta_n = n \lambda, \quad n = 0, \pm 1, \pm 2, \ldots,
\]  
где $d$ — расстояние между щелями, $\lambda$ — длина волны.  

Линейное расстояние $\delta x$ между соседними максимумами в фокальной плоскости линзы $O_2$ равно:  
\[
\delta x = \frac{\lambda f_2}{d},
\]  
где $f_2$ — фокусное расстояние линзы.  

\subsubsection*{Условия наблюдения}  
Чёткая интерференционная картина возникает только при выполнении условия когерентности:  
\[
d \leq \rho_{\text{ког}} \approx \frac{\lambda}{b} f_1,
\]  
где $b$ — ширина входной щели $S_2$, $f_1$ — фокусное расстояние коллиматора. При нарушении этого условия полосы размываются, что позволяет оценить размер источника света. Этот принцип используется в звёздных интерферометрах для измерения угловых размеров звёзд.  

\subsubsection*{Количество светлых полос в главном максимуме}  
Ширина главного дифракционного максимума (между первыми минимумами) определяется шириной отдельной щели $D$:  
\[
\Delta \theta = \frac{2\lambda}{D}.
\]  
Количество темных интерференционных полос $N$, укладывающихся в пределах главного максимума:  
\[
N = \frac{2d}{D}.
\]  
Таким образом, уменьшение ширины щели $D$ увеличивает число наблюдаемых полос.  

\subsection{Описание экспериментальной установки}  
\begin{itemize}  
    \item \textbf{Источник света}: Ртутная лампа с зелёным светофильтром ($\lambda = 546.1$ нм).  
    \item \textbf{Коллиматор}: Фокусное расстояние $f_1 = 12.5$ см.  
    \item \textbf{Экран с двумя щелями}:  
        \begin{itemize}  
            \item Ширина каждой щели: $D = 0.245 \pm 0.010$ мм (измерено в пункте Г).  
            \item Расстояние между щелями: $d = 0.99 \pm 0.10$ мм (измерено в пункте Г).  
        \end{itemize}  
    \item \textbf{Объектив}: Фокусное расстояние $f_2 = 11.5$ см.  
    \item \textbf{Микроскоп}: Окулярная шкала с ценой деления $0.02$ мм.  
\end{itemize}  

\subsection{Измерения и подсчёт результатов}  

\begin{figure}[h]  
\centering  
\includegraphics[width=1\textwidth]{interference_pattern}  
\caption{Интерференционная картина дифракции Фраунгофера на двух щелях}  
\label{fig:interference_pattern}  
\end{figure}  

\subsubsection*{Ход работы}  
\begin{itemize}  
    \item \textbf{Определение числа интерференционных полос}:  
        В пределах центрального дифракционного максимума наблюдалось $N = 7$ светлых полос.  

    \item \textbf{Расстояние между минимумами}:  
        \begin{align*}  
            \Delta x &= 43 \, \text{дел.} \cdot 0.02 \, \text{мм/дел.} = (0.86 \pm 0.04) \, \text{мм}, \\  
            \delta x &= \frac{\Delta x}{N} = \frac{0.86}{7} \approx (0.123 \pm 0.006) \, \text{мм}.  
        \end{align*}  

    \item \textbf{Расчёт расстояния между щелями}:  
        Используя формулу для интерференционных минимумов:  
        \[  
            d = \frac{2 \lambda f_2}{\delta x} = \frac{2 \cdot 546.1 \times 10^{-9} \, \text{м} \cdot 0.115 \, \text{м}}{0.123 \times 10^{-3} \, \text{м}} \approx (1.02 \pm 0.05) \, \text{мм}.  
        \]  

    \item \textbf{Влияние ширины источника}:  
        \begin{itemize}  
            \item Первое исчезновение полос при $b_0 = 0.91$ мм.  
            \item Восстановление картины при $b_1 = 1.02$ мм.  
        \end{itemize}  
\end{itemize}  

\subsubsection*{Теоретический анализ}  
\begin{itemize}  
    \item \textbf{Проверка $b_0$}:  
        Теоретическая ширина щели для исчезновения полос:  
        \[  
            b_0^{\text{теор}} = \frac{\lambda f_1}{d} = \frac{546.1 \times 10^{-9} \, \text{м} \cdot 0.125 \, \text{м}}{0.9875 \times 10^{-3} \, \text{м}} \approx (0.069 \pm 0.003) \, \text{мм}.  
        \]  
        Экспериментальное значение: $b_0 = 0.91 \, \text{мм}$. 

    \item \textbf{Проверка $b_1$}:  
        Решая уравнение для восстановления контраста:  
        \[  
            \frac{\pi d}{f_1} \cdot \frac{b_1}{\lambda} = 4.49341 \implies b_1^{\text{теор}} = \frac{4.49341 \cdot \lambda f_1}{\pi d} \approx (0.099 \pm 0.005) \, \text{мм}.  
        \]  
        Экспериментальное значение: $b_1 = 1.02 \, \text{мм}$.  
    \item \textbf{Проверка $N$}:
        Теоретическое количество полос в главном максимуме:
        \[
            N = \frac{2d}{D} \approx (8.1 \pm 0.1).
        \]
\end{itemize}  


\begin{longtable}[c]{|c|c|c|}  
\caption{Результаты измерений}  
\label{tab:double_slit_data} \\  
\hline  
Параметр & Значение & Погрешность \\ \hline  
\endfirsthead  

\multicolumn{3}{c}{{\tablename} \thetable\ — продолжение} \\  
\hline  
Параметр & Значение & Погрешность \\ \hline  
\endhead  

Число полос $N$ & 7 & 1 \\ \hline  
$\Delta x$, мм & 0.86 & 0.04 \\ \hline  
$\delta x$, мм & 0.123 & 0.006 \\ \hline  
$d$, мм & 1.02 & 0.05 \\ \hline  
$b_0$, мм & 0.91 & 0.05 \\ \hline  
$b_1$, мм & 1.02 & 0.05 \\ \hline  
\end{longtable}  

 Итак,
\begin{itemize}  
    \item \textbf{Расчётное расстояние между щелями} ($d = 1.02$ мм) близко к измеренному ($d = 0.9875$ мм).  
    \item \textbf{Теоретические $b_0$ и $b_1$} ($0.069$ мм и $0.099$ мм) значительно меньше экспериментальных ($0.91$ мм и $1.02$ соответственно), что указывает на возможную ошибку при экспериментальном измерении b. В обоих случаях ошибка составила ровно на порядок, скорее всего, ширина щели была получена неверно при выполнении работы.
    \item \textbf{Теоретическое число светлых полос в главном максимуме} N = 8 совпадает с экспериментальным значением (7) с точностью до одной полосы.
\end{itemize}  

\section{Пункт Г: Влияние дифракции на разрешающую способность оптического инструмента}

Дифракционные явления на оптических элементах принципиально ограничивают разрешающую способность любых инструментов. Рассмотрим установку, где коллиматор $O_1$ с фокусным расстоянием $f_1 = 12.5$ см формирует параллельный пучок света от двух щелей экрана $S_1$, расположенных на расстоянии $d = 0.9875$ мм друг от друга (будет измерено в этом пункте). Объектив $O_2$ ($f_2 = 11.5$ см) фокусирует изображение в плоскости наблюдения, при этом перед ним расположена регулируемая щель $S_2$ шириной $b$, дифракция на которой искажает изображение.

\subsection{Описание экспериментальной установки}
\begin{figure}[h]
\centering
\includegraphics[width=1\textwidth]{experimental_setup}
\caption{Схема установки для исследования разрешающей способности: 
1 — коллиматор $O_1$, 2 — экран $S_1$ с двумя щелями, 
3 — регулируемая щель $S_2$, 4 — объектив $O_2$, 
5 — плоскость наблюдения.}
\label{fig:setup}
\end{figure}

Основные элементы установки (рис. \ref{fig:setup}):
\begin{itemize}
\item Коллиматор $O_1$ ($f_1 = 12.5$ см) с экраном $S_1$, содержащим две щели ($d = 0.9875$ мм).
\item Регулируемая щель $S_2$ с микрометрическим винтом (ширина $b$ изменяется от $0.02$ до $1.5$ мм).
\item Объектив $O_2$ ($f_2 = 11.5$ см), формирующий изображение в фокальной плоскости.
\item Микроскоп с окулярной шкалой (цена деления $0.02$ мм) для анализа изображения.
\end{itemize}

Минимальное угловое расстояние $\varphi$ между точками объекта, которые еще воспринимаются как раздельные, определяется геометрией системы:
\[
\varphi \approx \frac{d}{f_1}.
\]
В фокальной плоскости объектива это соответствует линейному расстоянию между центрами изображений:
\[
l = f_2 \cdot \varphi = \frac{f_2 d}{f_1}.
\]

Угловая полуширина центрального дифракционного максимума от щели $S_2$ составляет:
\[
\Delta \varphi \sim \frac{\lambda}{b},
\]
где $\lambda = 546.1$ нм - длина волны света. Согласно критерию Рэлея, система способна разрешить две щели, если выполняется условие:
\[
\frac{\lambda}{b} < \frac{d}{f_1}.
\]
При нарушении этого неравенства изображения сливаются, делая невозможным различение двух щелей. Теоретический предел разрешения достигается при равенстве угловых размеров дифракционного пятна и расстояния между объектами, что соответствует критической ширине щели:
\[
b_{\text{крит}} = \frac{\lambda f_1}{d}.
\]

Таким образом, уменьшение ширины $b$ улучшает разрешающую способность системы, но одновременно усиливает дифракционные искажения. Этот фундаментальный компромисс между разрешением и искажениями является ключевым при проектировании любых оптических систем.

\subsection{Экспериментальные измерения}  
Для исследования разрешающей способности оптической системы были проведены следующие измерения:  
\begin{itemize}  
    \item Минимальная ширина щели $S_2$, при которой изображения двух щелей ещё различимы:  
    \[  
    b_0 = 1.31 \, \text{мм}.  
    \]  
    \item Координаты границ щелей, полученные с помощью микрометрического винта микроскопа:  
    \begin{align*}  
        \text{Левая щель:} & \quad 0 \, \text{мм} \, (\text{левая граница}), \, 0.245 \, \text{мм} \, (\text{правая граница}), \\  
        \text{Правая щель:} & \quad 0.99 \, \text{мм} \, (\text{левая граница}), \, 1.23 \, \text{мм} \, (\text{правая граница}).  
    \end{align*}  
\end{itemize}  

\subsection{Обработка результатов}  
\begin{itemize}  
    \item \textbf{Ширина каждой щели}:  
    \begin{align*}  
        D_{\text{лев}} &= 0.245 \, \text{мм} - 0 \, \text{мм} = 0.245 \, \text{мм}, \\  
        D_{\text{прав}} &= 1.23 \, \text{мм} - 0.99 \, \text{мм} = 0.24 \, \text{мм}.  
    \end{align*}  
    Средняя ширина щели:  
    \[  
    D = \frac{0.245 + 0.24}{2} = 0.242 \, \text{мм}.  
    \]  

    \item \textbf{Расстояние между щелями}:  
    \[  
    d = \frac{1.23+0.99}{2} - \frac{0245 + 0}{2}= 0.99 \pm 0.10 \, \text{мм}.  
    \]  

    \item \textbf{Проверка критерия Рэлея}:  
    Теоретическая критическая ширина щели $S_2$:  
    \[  
    b_{\text{крит}} = \frac{\lambda f_1}{d} = \frac{546.1 \times 10^{-9} \, \text{м} \cdot 0.125 \, \text{м}}{0.99 \times 10^{-3} \, \text{м}} \approx (0.069 \pm 0.003) \, \text{мм}.  
    \]  
    Экспериментальное значение $b_0 = 1.31 \, \text{мм}$ значительно превышает теоретический предел, что снова указывает на некорректность измерения ширины входной щели (как и в предыдущем пункте, ошибка произошла на порядок).  
\end{itemize}  

\begin{figure}[h]  
\centering  
\includegraphics[width=1\textwidth]{slit_measurement}  
\caption{Схема измерения координат щелей}  
\label{fig:slit_measurement}  
\end{figure}  